\documentclass[11pt,showkeys]{essay}

\begin{document}

\title{\textit{Human-Centered Theater:}\\
Overcoming Organizational Resistance to Authentic UX}

\author{Jaxen Dutta\, \orcidlink{0009-0002-5088-4679}}
\email[]{adutt042@uOttawa.ca}
\affiliation{School of Engineering Design and Teaching Innovation,\\ 
University of Ottawa, Ottawa, Ontario, Canada}

\date{September 21, 2025}

\begin{abstract}
Despite twenty-five years of advocacy for human-centered design principles, organizational adoption remains paradoxically limited across industries. This essay challenges conventional narratives that attribute slow implementation to organizational ignorance. Instead, it proposes that resistance reflects rational calculations about cultural transformation costs and the disruption of entrenched power structures. Through a systematic examination of design thinking's benefits, adoption barriers, and implementation strategies, this critique argues that authentic human-centered approaches face resistance due to their intensive implementation requirements rather than methodological shortcomings. While superficial ``human-centered theater'' allows organizations to claim user-centricity without meaningful change, authentic implementation requires fundamental shifts in organizational values and measurement systems that many leadership teams are unwilling to undertake.
\end{abstract}

\keywords{human-centered design, user experience, design thinking, organizational resistance, change management, planned obsolescence}

\maketitle

\section{Introduction}

Organizational resistance to \gls{hcd} reflects a systematic prioritization of operational efficiency over transformational \gls{ux} methodologies, even when the evidence for design thinking's impact remains overwhelming. Research by Hogan et al. reported return-on-investment (ROI)figures as high as 9,900\% for \gls{ux} investments, though these figures represent best-case scenarios under optimal conditions~\cite{hogan2016}. Nevertheless, organizations consistently resist the cultural changes required for authentic \gls{ux} implementation while simultaneously lamenting products that fail to meet user needs and demand costly redesigns.

When implemented authentically, \gls{hcd} delivers extraordinary competitive advantages through systematic approaches to understanding user behaviour, needs, and goals~\cite{illingworth2025wk1}. As a structured process, authentic \gls{ux} creates reliability across teams, provides scaffolding for continuous learning, establishes shared conceptions of project goals, and enables measurement capabilities for both throughput and quality~\cite{illingworth2025wk2}. It is a deliberate approach that encompasses all tangible and intangible aspects of user interaction.

\section{The Pitfalls of Human-Centered Theater}

Widespread organizational resistance reveals systematic implementation failures rather than limitations in the methodology itself. Most problematically, \gls{hcd} often degenerates into ``human-centered theater.'' These elaborate performances of user focus mask a fundamental disconnection from actual user needs. Corporate teams conduct extensive user research not to genuinely understand requirements, but to validate predetermined solutions and satisfy process requirements imposed by leadership teams that fundamentally misunderstand the cultural demands of user-centered design.

This emphasis on performance over authentic outcomes prioritizes procedural compliance, leading to situations where emotional language concerning desires and aspirations becomes a substitute for concrete problem-solving~\cite{kolko2015}. This phenomenon represents an insidious threat because it allows organizations to claim user-centricity while systematically ignoring users, fostering cynicism among stakeholders. Organizations frequently mistake narrow usability concerns for comprehensive \gls{ux} design, implementing expensive cosmetic improvements that fail to address fundamental product-market fit. 

Furthermore, the resource intensity argument against \gls{hcd} often masks a deeper unwillingness to confront necessary cultural changes. Design thinking requires a fundamental transformation involving organizational structure, processes, and mindsets~\cite{elsbach2018}. Yet, most organizations approach it as a mere process addition. The cross-functional collaboration required for authentic design thinking eliminates the massive inefficiencies created by siloed development processes, but it also challenges established departmental boundaries and power structures.

\section{Systematic Resistance to Design Evangelism}

Beyond internal organizational dynamics, the most fundamental barrier to \gls{ux} adoption lies in entrenched business models that directly contradict user-centered principles. Planned obsolescence represents perhaps the most egregious example. This practice deliberately designs products with limited lifespans to encourage frequent replacement purchases, prioritizing manufacturer revenue over user needs~\cite{slade2006}. 

Contemporary examples reveal this resistance across industries. Consumer electronics frequently adopt non-repairable designs through soldered components and proprietary tooling~\cite{perzanowski2021}. Automotive companies require dealer-specific diagnostic tools that contradict user preferences for accessible maintenance~\cite{europarl2024}. Software companies implement forced obsolescence through compatibility breaks that prioritize recurring revenue over user productivity. Organizations built around these models resist \gls{ux} adoption because authentic user research would reveal preferences that conflict with their revenue dependencies.

Despite this, successful implementations demonstrate that authentic user-centered design delivers measurable competitive advantages when conditions align. Design thinking helps organizations overcome cognitive biases that typically impede innovation~\cite{liedtka2018}. Notable examples include GE Healthcare's 90\% improvement in pediatric Magnetic Resonance Imaging (MRI) patient satisfaction~\cite{brown2008}, IBM's 75\% reduction in development time~\cite{datar2021}, Apple's market leadership via emotional connection~\cite{thomke2009}, and Netflix's transition to streaming driven by personalized \gls{ux}~\cite{gulati2019}. However, these successes required comprehensive cultural transformation initiatives.

\section{Recommendations for Authentic Implementation}

Presenting \gls{hcd} to organizational leadership requires rejecting several counterproductive accommodation strategies. Organizations should avoid translating design impact purely into traditional business language or minimizing the cultural change required. This perpetuates the dangerous illusion that design thinking provides equivalent results through different processes without fundamental transformation. 

Authentic implementation requires acknowledging that \gls{hcd} is a distinct business strategy. It trades short-term efficiency for long-term market position through superior user relationships and systematic advocacy. An assessment of organizational readiness is critical. This involves evaluating executive willingness to act on research findings that conflict with established assumptions, the capacity for sustained investment, and the cultural tolerance for challenging established hierarchies.

\section{Conclusion}

Persistent organizational resistance to authentic \gls{hcd} reflects a systematic avoidance of approaches that challenge existing power structures. The accommodation strategy of superficial ``human-centered theater'' only damages the long-term prospects for user-centered approaches. Success requires an honest assessment of organizational readiness and a sustained commitment to profound cultural transformation. Organizations unwilling to make these commitments should focus on operational excellence rather than attempting design implementations that are predisposed to fail.

\bibliography{human-centered-theater}
\end{document}
